\section{Introduction}

Large language models (LLMs) are powerful in storing and utilizing knowledge, yet their knowledge needs to continuously evolve with the changing world. This has led to widespread use of post-training techniques such as \weibing{fine-tuning and knowledge editing \cite{meng2022locating, meng2022mass, mitchell2021fast, de2021editing}}[todo: Added ROME, MEMIT, MEND citations] \heng{add citations} to update model knowledge. However, updating a single piece of knowledge often introduces unintended side effects, where changes propagate beyond the target fact and distort other, seemingly unrelated knowledge. \kmnote{This is the continual learning problem and appropriate references should be inserted. }This phenomenon, commonly referred to as the \textit{ripple effect}, highlights a fundamental challenge: knowledge updating in LLMs is not a local operation, but can trigger global changes in an interconnected system.

While ripple effects have been observed, it remains unclear \textit{what governs their propagation across knowledge}. \weibing{One natural hypothesis, often invoked informally, is that such effects are driven by surface-form and entity-name similarity}[todo: Reframed semantic similarity as one hypothesis rather than established fact, and clarified semantic to surface-form] between facts. Prior work and empirical observations suggest that models are prone to confusion among \weibing{surface-similar or name-similar entities}[todo: Avoided ambiguity with semantic similarity], leading to the expectation that errors may propagate along similarity (\textcolor{blue}{cite}). At the same time, another plausible hypothesis is that different types of knowledge exhibit different levels of robustness. In particular, \weibing{by analogy with the long-tail knowledge literature, one might expect graph-structural tails (i.e., entities with low in-degree within a given domain) to be more fragile}[todo: Softened long-tail assumption and defined it structurally] and thus more susceptible to corruption (\textcolor{blue}{cite}). \heng{these descriptions are all too abstract}
In this work, we aim to test these hypotheses and answer a central question:
\kmnote{The above section is nicely written so long as there are ties to continual learning. The difference may be in looking at relatedness in language. I don't think semantic relatedness is generally considered, though we should check. }

\textbf{What determines the ripple effect in knowledge updating?}

To systematically study how relationships between knowledge affect ripple propagation, we construct a structured knowledge graph grounded in real-world corpora. \weibing{To track exactly how an edit disrupts related facts, we inject controlled \textit{counterfactual} facts (plausible but incorrect updates) to simulate update errors and trace their propagation.}[todo: Defined counterfactual update in intro] This allows us to explicitly model connections between facts and quantify how updates propagate along these connections. Within this framework, each piece of knowledge corresponds to a node, and relationships between facts define the edges, enabling us to analyze how ripple effects unfold along the graph structure.

Using this framework, we analyze how ripple effects vary across different parts of the graph. 
We observe that \weibing{highly connected object nodes (Hubs)}[todo: Standardized to Hubs] are significantly more vulnerable to corruption than sparsely connected ones. \kmnote{Can you introduce the term "popularity" here? You might also more finally define, particularly so that readers can determine it is not frequency. You do mention this in abstract. it should occur early in intro}
Moreover, \weibing{Hubs}[todo: Terminology] not only exhibit greater susceptibility, but also play a dominant role in propagating errors to other parts of the knowledge space.

Interestingly, the connectivity of a node in the graph aligns with how frequently a piece of knowledge is referenced or shared across facts, allowing us to interpret it as a proxy for \weibing{graph-structural popularity}. \kmnote{OK I see it here. But use italics and more formally define. I think you do want to distinguish from frequency. I thought in Bengal talk you said it was not frequency. }
Under this interpretation, \weibing{Hubs}[todo: Terminology] correspond to widely shared or commonly used knowledge, while sparsely connected nodes correspond to \weibing{low-connectivity tails (Tails)}[todo: Terminology].
From this perspective, the observed behavior is counterintuitive: \weibing{structurally central}[todo: Terminology] knowledge is generally expected to be more robust, \heng{not sure what you mean by knowledge is robust} yet we find it to be more vulnerable and more influential in error propagation. 
This finding contradicts both the similarity-based explanation \kmnote{what is this? explain and cite. }and the common intuition that \weibing{tail knowledge}[todo: Terminology] is the primary source of fragility, \kmnote{cite this about long-tailed also. }suggesting instead that ripple effects are governed by how centrally knowledge is connected. 
This observation naturally leads to a deeper question:

\textbf{Why is \weibing{highly connected knowledge} both more vulnerable and more influential?}


To answer this, we analyze the internal representations of LLMs from two complementary perspectives. First, we show that \weibing{Hubs} \kmnote{What is hub? You are using many different terms. Is this popular? Choose one term and stick with it. }exhibit significantly narrower decision margins, indicating fragile decision boundaries that are easily perturbed by updates. \kmnote{This previous sentence not clear. What are narrower decision margins? }Second, \weibing{as exploratory mechanistic evidence, we examine whether attention perturbations align with the observed propagation patterns, suggesting pathways}[todo: Softened attention claim] for propagating updates across distant knowledge. Together, these results suggest that the same structural property, connectivity, simultaneously determines both the stability of knowledge representations and the pathways through which errors spread.
Finally, this understanding raises a practical question:


\textbf{Can we leverage \weibing{Hubs} to control ripple effects?}


If highly connected nodes amplify error propagation, they may also serve as anchors to stabilize the model. Building on this insight, we propose a connectivity-aware regularization strategy that selectively anchors \weibing{Hubs} during updates. We show that this approach effectively reduces long-range ripple effects and improves overall stability, outperforming connectivity-agnostic baselines. This demonstrates that the same mechanism that causes instability can also be harnessed for control.

Overall, our work reveals a unifying principle:

\textbf{Node popularity, as a proxy for structural connectivity, governs both the vulnerability and propagation of ripple effects in LLMs.}

By identifying connectivity as the key factor underlying knowledge distortion, we provide both a deeper understanding of knowledge dynamics in LLMs and a practical pathway toward more reliable knowledge updating.

\kmnote{Overall this intro is very interesting. Just needs tightening up of terminology use, more citations, and definitions.}